\chapter*{Resumen}
\addcontentsline{toc}{chapter}{Resumen}

Este Trabajo Fin de Máster aborda el estudio práctico de las vulnerabilidades
web más comunes mediante el diseño, la implementación y la explotación
controlada de un laboratorio propio, denominado \breachlab. La aplicación,
compuesta por un backend en Flask y un frontend en Astro servidos bajo el
mismo origen mediante nginx, implementa de forma deliberada seis fallos de
seguridad: inyección SQL, cross-site scripting (XSS), cross-site request
forgery (CSRF), fallos de autenticación, control de acceso roto (IDOR) y una
configuración de seguridad insegura; junto con su corrección correspondiente,
activable mediante un parámetro dentro del enlace \texttt{(?safe=1)}, lo que permite al 
usuario comparar de forma directa el comportamiento vulnerable y el seguro sobre el mismo código.

A partir del estudio de la OWASP Web Security Testing Guide (WSTG), tomada
como metodología de referencia, se deriva una metodología propia más ágil:
un ciclo fijo de seis pasos --- reconocimiento, identificación, explotación,
evidencia, verificación de la mitigación y valoración de impacto ---
aplicado a un subconjunto acotado de pruebas alineado con el OWASP Top 10,
pensado para equipos de desarrollo pequeños o pymes con recursos limitados.
Para cada vulnerabilidad se registran evidencias objetivas mediante un
sistema de logging propio y se calcula su severidad conforme al estándar
CVSS v3.1. Los resultados se comparan, por un lado, con la propia OWASP
WSTG completa, en cobertura y esfuerzo estimado, y por otro con dos
laboratorios de referencia del sector, DVWA y OWASP Juice Shop; y se
ilustran mediante un caso de estudio que encadena varias de las
vulnerabilidades analizadas para reproducir un escenario de compromiso
realista. El trabajo concluye con un conjunto de recomendaciones de
mitigación y propuestas de ampliación del laboratorio y de la metodología.

\bigskip
\noindent\textbf{Palabras clave:} seguridad web, metodología ágil,
OWASP WSTG, OWASP Top 10, inyección SQL, XSS, CSRF, IDOR, CVSS, laboratorio
de pruebas, seguridad ofensiva.
